\chapter{Chemistry vs Physics}\label{ch:chemvphys}
\index{chemistry vs physics}\index{Dirac 1929 declaration}\index{applied physics, chemistry as}

\begin{quote}
\itshape It seems to me that the concept of probability is terribly mishandled these days. A probabilistic assertion presupposes the full reality of its subject. No reasonable person would express a conjecture as to whether Caesar rolled a five with his dice at the Rubicon. But the quantum mechanics people sometimes act as if probabilistic statements were to be applied \emph{just} to events whose reality is vague.\\
\normalfont\hfill --- E.~Schr\"odinger to A.~Einstein, 1950.
\end{quote}
\bigskip

A central question in the philosophy of chemistry is whether \textbf{(A)} chemistry can be explained by quantum physics, as Dirac famously claimed in 1929 at the birth of quantum mechanics, or \textbf{(B)} something more is required, as has been argued by leading chemists for nearly a century.

Position (B) is the view that chemistry possesses an irreducible conceptual content --- orbitals, bonds, valence, hybridisation, electronegativity --- that does not follow from Schr\"odinger's equation alone, no matter how many computational resources are thrown at it. Position (A) is the view that chemistry \emph{is} applied quantum physics, and that the apparent autonomy of chemical concepts is an artefact of the practical impossibility of solving the underlying equations.

The hundred-year stalemate between (A) and (B) has been driven less by disagreement on principle than by the fact that, within the standard $3N$-dimensional wave-function formulation that Dirac himself had in mind in 1929, the equations are indeed unsolvable for any system of chemical interest without approximations whose chemistry-specific character (basis sets, exchange--correlation functionals, fitted force fields) seems to vindicate position (B).

This chapter argues that RealQM, by reformulating the problem in real 3D space with non-overlapping unit densities, provides concrete evidence that Dirac's position (A) was correct --- although not within the 1929 wave-function-on-$3N$-space formulation that Dirac himself proposed. Chemistry emerges from physics; the equations were just being written in the wrong space.

\section{The two positions}\label{sec:two-positions}

The chemistry-versus-physics debate has been long-running, often heated, and rarely settled. The clearest articulations of the two positions are themselves widely quoted, and we restate them here to anchor the discussion.

\paragraph{Position (A): chemistry as applied physics (Dirac, 1929).}
\begin{quote}
The underlying physical laws necessary for the mathematical theory of [\ldots] the whole of chemistry are thus completely known, [and] the difficulty is only that the exact application of these laws leads to equations much too complicated to be soluble.
\end{quote}
Dirac's claim is unambiguous: the laws are known (the Schr\"odinger equation and Coulomb interactions); the obstacle is computational. Chemistry is what one gets when one applies known physical laws to many-electron systems.

\paragraph{Position (B): chemistry has irreducible content (the chemists' counter).}
A typical articulation of position (B), drawn from the long line of chemical-philosophy literature: orbitals, bonds, valence, hybridisation, and electronegativity are concepts of chemistry that are not derivable from the Schr\"odinger equation. They are not artefacts; they are the working concepts that organise chemical knowledge, and they do work that the underlying physics does not do. A reduction of chemistry to physics is not on offer in the conventional sense; chemistry is autonomous from physics, even if it is consistent with physics.

The proponents of position (B) have been distinguished chemists, including R.~G.~Woolley, H.~Primas, J.~Earley, and many others. The argument is not anti-physics; it is that chemistry has a conceptual structure that emerges at a higher level than the underlying physics and is not derivable downward.

\paragraph{Schr\"odinger's persistent objection.}
Schr\"odinger himself, who wrote the equation in 1926, never accepted the probabilistic interpretation that Born and Heisenberg attached to it. The opening epigraph of this chapter is one such statement from 1950; the chapter on RealQM vs Statistics collects further passages from his 1935 paper, and additional quotations open Chapter~\ref{ch:nucleus} (Wien letter 1926) and Chapter~\ref{ch:collab} (Bohr 1926). Position (B) above is, in the end, an institutional response to a foundation that the founder himself never endorsed. RealQM's contribution to the debate is to provide a working alternative in the formulation Schr\"odinger had argued for.

\section{Why the debate persisted}\label{sec:why-persisted}

The debate persisted for nearly a century because, within the standard formulation of quantum mechanics, the Schr\"odinger equation on $\mathbb{R}^{3N}$ \emph{is} unsolvable for systems of chemical interest. The exponential scaling of configuration space puts any system beyond a few electrons out of reach for exact treatment. The methods of standard quantum chemistry --- Hartree--Fock, density functional theory, configuration interaction, coupled cluster --- are approximations that introduce chemistry-specific machinery to make the calculations tractable:

\begin{itemize}
\item Basis sets (Gaussian, plane-wave, Slater-type) are chemistry-specific choices.
\item Exchange--correlation functionals in DFT are fitted to chemistry-specific data.
\item Pseudopotentials and effective core potentials encode chemistry-specific assumptions about which electrons participate in bonding.
\item Density-fitting auxiliary basis sets, dispersion corrections, range separation --- all chemistry-specific.
\end{itemize}

When the working methods of quantum chemistry are themselves saturated with chemistry-specific choices, the claim that chemistry reduces to physics becomes hard to support. Position (B) gains plausibility every time a quantum chemist makes a choice about basis sets or functionals: the chemistry-specific knowledge is doing real work in the calculation, and the calculation would not produce useful results without it.

The standard rejoinder from position (A) is that the chemistry-specific machinery is computational, not physical: the physics is the Schr\"odinger equation, and the approximations are practical compromises. But this rejoinder runs into the wall that Dirac himself acknowledged: \emph{the equations are too complicated to be soluble}. As long as the equations cannot be solved exactly without approximation, position (A) remains a statement of principle rather than a working demonstration.

\section{RealQM as evidence for (A) --- in a different formulation}\label{sec:realqm-evidence}

RealQM provides \textbf{concrete evidence that Dirac's position (A) was correct, although not within the 1929 wave-function-on-$3N$-space formulation that Dirac himself proposed}. The reformulation as a system of non-overlapping unit electron densities in real 3D space, with energy minimisation over Coulomb interactions alone, recovers chemistry from physics without chemistry-specific empirical input.

Specifically:

\begin{itemize}
\item Atomic energies (Chapter~\ref{ch:atoms}) follow from minimum-energy packing of non-overlapping unit densities around a positive kernel. The shell structure $2, 8, 18, 32$ and the sub-shell radial split $(4{+}4)$ are emergent predictions, not put in by hand. The reduction is at 1\% accuracy across Li--Rn with no fitted parameters.

\item The covalent bond and its Helium limit (Chapter~\ref{ch:h2}) follow from the variational principle when two nuclei are present. The Bernoulli condition at the inter-electronic free boundary --- which itself emerges from the variational principle --- makes the bond energetically favourable through density accumulation between the kernels. The H$_2$ result agrees with Kolos--Wolniewicz to 0.3\% with no fitted parameters.

\item Atomization energies of hydrides across periodic-table groups (Chapter~\ref{ch:hydrides}) follow from a Level-3 reduction in which the kernel softening radius $r_c$ traces a column of the periodic table. The architecture is per-group, not per-molecule, with the chemistry-specific machinery reduced to a single parameter per atomic species.

\item Reactive chemistry (Chapter~\ref{ch:dimers}) follows from the local-Coulomb force law of Chapter~\ref{ch:equation}. The same gradients that hold a bond together also break it and re-form it. No transition-state machinery, no pK$_a$ table, no reaction-specific input.

\item Protein folding (Chapter~\ref{ch:folding}) follows from quantum-mechanical H-bond forces on backbone atoms, with one implicit-solvent parameter (the SASA hydrophobic pressure) to compensate for the entropic water effect that the framework does not represent explicitly.

\item Condensed-phase physics (Chapter~\ref{ch:condensed}) follows from the same architecture with Brownian dynamics. Phase transitions emerge with the right Lindemann and RDF character; the overshoot factor against experimental transition temperatures quantifies the missing-dispersion contribution.

\item Stability of matter (Chapter~\ref{ch:som}) follows from Hardy's inequality applied to non-overlapping unit densities. The proof is one paragraph; the underlying physical content is the same as in the StdQM Dyson--Lenard--Lieb--Thirring proof, but the formulation is simpler.

\item The nuclear extension (Chapter~\ref{ch:nucleus}) shows that the same multiphase formulation applies, with proton--electron constituents, to atomic nuclei. The framework therefore reaches at least two scales of physics through a single variational principle.
\end{itemize}

Across this list, chemistry emerges from physics --- from the variational minimisation of the Coulomb energy on a non-overlapping partition of physical space --- with no chemistry-specific empirical input beyond the kernel parameters that encode the level-by-level reduction. The chemistry is not autonomous from the physics; it is what the physics produces when it is applied to multi-nucleus systems.

\section{What changed: the right space}\label{sec:right-space}

The key change is the space the wave function lives on. Dirac's 1929 statement was made in the context of the $3N$-dimensional formulation that he, Heisenberg, and others were developing at the time. Within that formulation, the underlying physics is indeed known and the equations are indeed too complicated to be soluble. The standard quantum chemistry that followed --- Hartree--Fock, configuration interaction, coupled cluster, density functional theory --- is the working compromise that emerges from trying to solve the $3N$-dimensional equation by truncation and approximation.

The multiphase formulation moves the wave function back to physical 3D space, in the spirit of Schr\"odinger's 1926 letter to Lorentz (Chapter~\ref{ch:dirac}). The wave function $\Psi(\mathbf{x}) = \sum_i \psi_i(\mathbf{x})$ on $\mathbb{R}^3$ has the dimensionality of physical space; the cost of representing it on a grid is $O(K)$ rather than $O(K^N)$. The equations are no longer too complicated to be soluble: the entire framework fits in 8000 lines of code and runs interactively on a laptop GPU.

\textbf{Dirac was right; the equations were just being written in the wrong space.} The chemistry-specific machinery that StdQM had to introduce to make the $3N$-dimensional equations tractable --- the basis sets, the functionals, the pseudopotentials --- is absent in RealQM not because RealQM is more accurate (it is not, in absolute terms) but because the reformulation does not need that machinery to be tractable. The chemistry comes out without it.

\section{What this does and does not say about chemistry's autonomy}\label{sec:autonomy}

The argument is more nuanced than a simple ``physics wins, chemistry is just applied physics.'' Let us be careful about what is and is not being said.

\paragraph{What is said.}
Chemistry-as-applied-physics is achievable when one is willing to reformulate the problem in real 3D space with an appropriate non-overlap constraint. The chemistry-specific machinery of StdQM is not necessary for the bulk of chemistry to emerge from physics; it is necessary for the bulk of chemistry to emerge from the \emph{$3N$-dimensional formulation} of physics. Choosing a different formulation eliminates the need for the machinery.

\paragraph{What is not said.}
\begin{itemize}
\item The standard $3N$-dimensional formulation should be abandoned. For high-precision spectroscopy, transition probabilities, and excited-state dynamics, it remains the canonical tool. RealQM does not reach CCSD(T)/CBS precision and is not claimed to.
\item Chemists' working concepts (orbitals, bonds, hybridisation, electronegativity) are wrong. They are heuristic shorthand for emergent structures in the underlying many-electron problem, and they work because the structures they describe are real. RealQM does not deny their utility; it provides an alternative route to the same structures.
\item Chemistry has no autonomy of any kind. The hierarchy of reductions in Chapter~\ref{ch:hierarchy} is itself a layered description, with each level having its own working concepts (shell structure, pseudo-kernel, splitting topology). The reduction goes from physics to chemistry, but the chemistry that emerges is organised at its own level. Position (B) is partially correct: chemistry has its own working concepts. Position (A) is also correct: those concepts emerge from physics.
\end{itemize}

The two positions are not strictly incompatible. The disagreement is about whether the emergence is achievable in practice, and what is required for it to be achievable. Position (A) won as a matter of principle if the formulation can be chosen freely. Position (B) won as a matter of practice for as long as the only available formulation was the $3N$-dimensional one. RealQM shifts the balance by demonstrating that the practical obstacle was specific to the formulation, not to the underlying claim.

\section{Implications for chemistry education}\label{sec:education}

If the chemistry-versus-physics debate is resolved in the way this chapter argues, there are concrete implications for how chemistry is taught.

\paragraph{Stability of matter as introductory result.}
As discussed in Chapter~\ref{ch:som}, stability of matter is currently absent from standard QM curricula because the StdQM proof is too difficult. In RealQM, it is a one-paragraph consequence of Hardy's inequality. A first-year student can be shown why matter is stable.

\paragraph{The covalent bond from a free-boundary condition.}
The standard pedagogical account of the covalent bond invokes orbital overlap, bonding and antibonding molecular orbitals, exchange, and antisymmetry. In RealQM, the bond is what happens when two electron territories meet at a Bernoulli free boundary between two positive kernels: the non-zero meeting allows density to accumulate between the kernels at low kinetic-energy cost, lowering the total energy. The mechanism is direct and visualisable.

\paragraph{Periodic-table trends from kernel softening.}
The periodic-table trend within a group (lighter atoms bind harder, heavier atoms bind softer) is, in RealQM, a sweep of a single parameter $r_c$. The trend has a structural cause --- inner-shell absorption --- that can be made visible interactively in the browser-based implementation.

\paragraph{Chemistry as continuum mechanics.}
The most ambitious pedagogical claim is that chemistry can be taught as a special case of continuum mechanics on a partitioned domain, with the Bernoulli free-boundary condition as the closing condition. The mathematical machinery is standard variational calculus and PDE theory; the chemistry is what emerges when the partitioned-domain problem is solved for multi-kernel Coulomb systems.

Whether this pedagogical reorientation is taken up is a separate question; the present book argues that it is achievable in principle.

\section{Implications for physics education}\label{sec:physics-education}

The implications for physics education are sharper than the chemistry case, because the standard report from inside the discipline is that StdQM \emph{is not understood}. The most famous statement comes from Richard Feynman:

\begin{quote}
\itshape I think I can safely say that nobody understands quantum mechanics.\\
\normalfont\hfill --- R.~Feynman, \emph{The Character of Physical Law}, 1965.
\end{quote}

The remark is not a rhetorical flourish. It is a working assessment by a Nobel laureate, on the lecture record, and it has been repeated in similar form by Gell-Mann (``mysterious, confusing\ldots which none of us really understands''), by Mermin (``shut up and calculate''), and by a long line of others quoted in the Prologue's strangeness column. Standard quantum mechanics is taught, used, and applied with extraordinary empirical success, and not understood. The two facts coexist in the same physicist.

The educational consequence is severe. A student is told from the first lecture that the central object of the theory --- the wave function on $3N$-dimensional configuration space --- has no direct physical meaning, that intuition trained on everyday mechanics will mislead, that the right response to the question ``but what is actually happening?'' is to stop asking the question. Physical-mechanics intuition, the asset every student arrives with, is positioned as the enemy. The professor passes on what they were taught: an apparatus that works and an injunction against asking what it means. The student passes it on in turn.

\paragraph{What RealQM offers the student and the teacher.}
RealQM removes this educational gap because its central object \emph{has} direct physical meaning. The wave function lives on physical 3D space; $\psi_i^2$ is a literal charge density of electron $i$ on its region $\Omega_i$; the electrons interact through the Coulomb potential they jointly generate, with the partition self-consistently adjusted by energy minimisation. The framework is, in the technical and the everyday sense, \textbf{interacting bodies on a shared physical stage}: distinct objects, each with its own spatial extent, each perceiving the rest through the field the rest generates, the configuration determined by mutual energetic balance. A first-year student arriving with intuition trained on classical mechanics --- bodies that occupy space, interact through fields, settle into equilibrium configurations --- finds that intuition rewarded rather than punished. The intuition extends; it does not have to be unlearned.

The teacher is in the same position. The Bernoulli free-boundary condition, the Coulomb potential between charge densities, the variational minimisation, and the partitioning of space into territories are all standard items in the undergraduate curriculum in classical continuum mechanics. A teacher fluent in classical mechanics is fluent in RealQM after reading the present volume; no separate apparatus of complex Hilbert spaces, antisymmetric tensor products, projection postulates, or interpretational layers is required. The teacher can answer the student's questions in the same vocabulary the questions were asked in.

\paragraph{RealQM as physics education in the everyday sense.}
The deepest pedagogical claim of this book is that quantum mechanics, as far as Coulomb chemistry goes, can be taught as \textbf{everyday mechanics of interacting bodies in 3D space}. Electrons are bodies; they have extent; they interact through Coulomb fields they generate; the equilibrium configuration of a many-electron system is the energy-minimising partition. The hydrogen atom is a single charge cloud bound to a proton; the Helium atom is two such clouds on two complementary half-spaces; an L-shell atom is eight such clouds at the vertices of a Lewis cube; a covalent bond is what happens when two such clouds meet at a Bernoulli boundary between two protons; a melting solid is the same partition coming undone at a critical temperature. None of these statements requires interpretational machinery; all of them can be drawn on a blackboard, simulated in a browser, and varied parametrically by the student. The framework is, in a strict and practical sense, physics that can be taught to and understood by physics teachers and physics students.

Whether the pedagogical reorientation is taken up depends on whether the physics community is willing to revisit the 1926--1929 choice. The present volume demonstrates that the formulation works and that a student equipped with a browser and the public Gallery can develop genuine understanding of multi-electron chemistry as everyday continuum mechanics. A century of ``nobody understands quantum mechanics'' is a long time for the discipline to have taught a theory whose own founders thought could not be understood; that the alternative \emph{can} be understood, and recovers the same chemistry across the regimes that matter for physics and chemistry education, is the corollary that follows directly from the validation in Part~III.

\section{Hilbert spaces: 3D physical space vs.\ \texorpdfstring{$3N$}{3N}-dimensional configuration space}\label{sec:hilbert-distinction}

A standard claim about quantum mechanics is that it lives in Hilbert space, and that this distinguishes it structurally from classical continuum physics. The claim is, on inspection, misleading. The Hilbert-space machinery used in standard quantum mechanics divides cleanly into two parts of very different status, and the framework's reputation as conceptually exotic attaches to one of them and not the other.

\paragraph{Hilbert spaces of functions on physical 3D space.}
Standard PDE theory for classical continuum mechanics --- the Navier--Stokes equations on a fluid domain, the equations of linear elasticity on a deformable body, the heat equation on a thermally conductive region, the buckling eigenproblem of a plate --- naturally lives in Sobolev spaces $H^1(\Omega)$, $H^2(\Omega)$, and similar function spaces over a domain $\Omega \subset \mathbb{R}^3$. These are Hilbert spaces. The Lax--Milgram theorem, the Riesz representation theorem, the spectral theorem for self-adjoint elliptic operators, and the variational formulation of eigenvalue problems are all theorems about these Hilbert spaces and apply uniformly across PDE-based continuum mechanics. The Schr\"odinger eigenvalue problem on physical 3D space --- the hydrogen problem $-\tfrac{1}{2}\Delta\psi - \psi/|\mathbf{x}| = E\psi$, the RealQM functional~\eqref{eq:TE} minimised over $\psi_i \in H^1(\Omega_i)$ --- is an instance of this same standard PDE-Hilbert-space machinery. Nothing about it is structurally distinct from the eigenvalue theory of classical elastic plates or fluid buckling problems. It is variational calculus of a Sobolev-space functional, solved by the same techniques and analysed with the same theorems.

\paragraph{Hilbert spaces over \texorpdfstring{$3N$}{3N}-dimensional configuration space.}
The Hilbert space invoked in standard quantum mechanics is something else: the space $L^2(\mathbb{R}^{3N})$ of square-integrable functions of $3N$ coordinates, with an inner product that pairs functions on this abstract $3N$-dimensional configuration space. This is a Hilbert space in the formal sense --- inner product, completeness, the same axioms --- but it is a fundamentally different object from a Sobolev space of fields on physical 3D space. Its elements are not fields on physical space; they are scalar-valued functions of a $3N$-dimensional abstract coordinate. The probabilistic interpretation of $|\Psi|^2$ as a joint configuration density attaches to this Hilbert space, not to the Hilbert spaces of classical PDE theory. The non-commuting-observables apparatus, the projection postulate, the superposition principle, and the entire ``mysteries of quantum mechanics'' attach here. The exponential computational cost of representing $\Psi$ on a grid also attaches here.

\paragraph{The distinction matters.}
The Hilbert spaces of standard PDE continuum mechanics are spaces of fields on physical 3D space. The Hilbert spaces of standard quantum mechanics are spaces of functions on $3N$-dimensional configuration space. Both are technically Hilbert spaces in the functional-analytic sense; in their physical content and computational character they are radically different. The reputation of QM as ``a theory in Hilbert space'' has been read as if it meant ``in the same Hilbert-space framework as classical continuum mechanics, only stranger,'' when what it actually means is ``in a different and far larger space whose dimensionality grows with particle count.'' The strangeness of QM is not a strangeness of Hilbert spaces in general; it is the consequence of the specific choice to extend the wave function from $\mathbb{R}^3$ to $\mathbb{R}^{3N}$.

\paragraph{Hilbert spaces have little to do with atomic physics as such.}
The conclusion is sharper than the preceding paragraphs may suggest. \textbf{Hilbert spaces in the standard PDE sense --- function spaces over physical 3D space --- have little to do with atomic physics specifically; they are the universal mathematical setting for variational continuum problems, and atomic physics in its real-space formulation is one continuum problem among many.} What \emph{is} specific to standard quantum mechanics is the extension to $3N$-dimensional configuration space, with the probabilistic interpretation that attaches to it. RealQM removes that extension and works in the standard PDE-Hilbert-space framework throughout. Atomic physics, in this formulation, sits next to fluid mechanics and elasticity in its mathematical kinship --- not as an isolated discipline with its own exotic functional-analytic apparatus, but as a particular continuum problem whose closing condition (the Bernoulli free boundary) is what makes it specifically atomic. The ``Hilbert space mystique'' of standard quantum mechanics dissolves on this view; what remains is variational calculus on Sobolev spaces and the question of how the partition of physical space closes.

\paragraph{An empty incantation in the lecture hall.}
The standard pedagogical incantation --- ``the wave function $\Psi$ lives in a Hilbert space'' --- deserves a closer reading than it usually receives. Phrased this way the sentence sounds impressive and specific; in fact it is neither. ``Lives in a Hilbert space'' is true of essentially every well-posed function in mathematical physics: the velocity field of a fluid, the displacement field of an elastic body, the temperature field of a heat-conduction problem, the eigenfunctions of any self-adjoint elliptic operator on a bounded domain, the solutions of every variational problem in continuum mechanics. The statement carries no information that distinguishes atomic physics from the rest of mathematical physics, and no information about what $\Psi$ actually represents. It is impressive-sounding terminology without physical content --- an empty incantation that the speaker often does not unpack and the audience often does not press to have unpacked. The genuine technical content of standard QM lies not in the words ``Hilbert space'' but in the silent qualifier $L^2(\mathbb{R}^{3N})$ that the incantation does not say out loud: the wave function lives on a $3N$-dimensional abstract space, with all the interpretational and computational consequences that follow. The lecture-hall phrasing is, in this strict sense, sloppy speech --- impressive in tone, evasive in content. Saying instead ``$\Psi$ is a scalar-valued function of $3N$ coordinates, interpreted as a joint probability density'' would convey the same mathematical content and make the strangeness audible from the first sentence.

\section{The von Neumann split: 1932 formalism vs.\ 1945 architecture}\label{sec:von-neumann-split}\index{von Neumann, J.}\index{von Neumann architecture}\index{EDVAC}

The case made in the previous section --- that the Hilbert-space framework attributed to QM is doubly removed from its claimed source --- is sharpened by following the trajectory of the man who actually built that framework. John von Neumann's career is, in retrospect, a tight illustration of the position this book takes.

\paragraph{1927--1932: the QM formalism.}
A short sequence of papers culminating in \emph{Mathematische Grundlagen der Quantenmechanik}~\cite{vonNeumann1932} established the operator-on-Hilbert-space formulation of quantum mechanics that has been standard ever since: states as unit vectors in $L^2(\mathbb{R}^{3N})$, observables as self-adjoint operators, measurement as orthogonal projection, the spectral theorem as the mechanism connecting operators to possible measurement outcomes. The book also contained the no-hidden-variables theorem (later contested by Bell). It was the definitive statement of the abstract framework into which the 1925--1929 work of Heisenberg, Born, Jordan, Schr\"odinger, and Dirac was consolidated.

\paragraph{1935: a private doubt.}
Three years after publishing the \emph{Grundlagen}, von Neumann wrote to Garrett Birkhoff a remark that has become one of the more striking documents of internal scepticism in the history of quantum mechanics~\cite{Redei2005}:

\begin{quote}
\itshape I would like to make a confession which may seem immoral: I do not believe absolutely in Hilbert space any more.\\
\hfill --- J.~von Neumann to G.~Birkhoff, 1935.
\end{quote}

This is the author of the operator-Hilbert-space formulation expressing doubt, in private correspondence, about the framework that bears his name in QM. The remark is not a casual aside: it is followed in von Neumann's subsequent work by a sustained move into alternative mathematical structures (continuous geometries, operator algebras as ``rings of operators'', quantum logic with Birkhoff in 1936) that are technically and conceptually different from the Hilbert-space framework of the \emph{Grundlagen}.

\paragraph{1936 onwards: the move into computing and applied mathematics.}
After the mid-1930s von Neumann's principal work was no longer in QM foundations. Game theory (\emph{Theory of Games and Economic Behavior}, with Morgenstern, 1944); operator algebras with F.~J.~Murray; the Manhattan Project (implosion-design hydrodynamics, 1943--1945); the EDVAC report~\cite{vonNeumann1945} establishing the architecture for stored-program computers; theory of automata; cellular automata; self-reproducing machines; numerical analysis (Monte Carlo, hyperbolic-PDE schemes); the first computer-based weather forecasts (with Charney) --- the bulk of his output for the next two decades is in fields his QM framework was not built to address. He became, in the second half of his career, the founding figure of \emph{computable} mathematics and \emph{computable} physics.

\paragraph{The structural incompatibility.}
The split is not a contingent fact of one man's biography; it is a structural problem inside the formalism itself. The 1932 Hilbert-space formulation makes the central object of QM --- the wave function $\Psi$ on $\mathbb{R}^{3N}$ --- cost $O(K^N)$ to represent on a $K$-point grid. The 1945 von Neumann architecture for computers --- finite memory, sequential instruction execution, real-time deterministic state evolution --- is the practical embodiment of computable mathematics. The same person built both, and the two pieces of his work are structurally incompatible for $N>1$ in the precise sense that the central object of the first cannot be represented on the second. Walter Kohn's 1998 line that $\Psi$ is ``not a legitimate scientific concept'' for $N \geq 10^2$--$10^3$ (Prologue, \S\ref{sec:prologue-strange}) is exactly the statement that von Neumann's 1932 QM is not runnable on von Neumann's 1945 computer.

\paragraph{RealQM as the reconciliation.}
The framework developed in this book is, in this lineage, the reconciliation von Neumann did not finish. By keeping the wave function on physical $\mathbb{R}^3$ as a sum of one-electron unit-density wave functions on a partition of space, with the variational principle closing through a free-boundary (Bernoulli) condition, RealQM produces a formulation of quantum mechanics that runs on actual von Neumann hardware --- a laptop GPU, with the 8000-line implementation described in the Appendix --- at interactive rates. The Hilbert spaces involved are the standard PDE-theory Hilbert spaces $H^1(\Omega_i)$ over physical 3D space, the same kind that organise classical continuum mechanics (Section~\ref{sec:hilbert-distinction}); the abstract $L^2(\mathbb{R}^{3N})$ Hilbert space that von Neumann himself doubted in 1935 is not invoked at all. The two halves of von Neumann's intellectual legacy --- the formalism and the architecture --- are made compatible by going back to Schr\"odinger's 1926 real-space picture, rather than continuing on the 1932 abstract Hilbert-space line that von Neumann himself had reservations about within three years of publishing.

\section{Mathematicians and StdQM: the engagement that did not happen}\label{sec:mathematicians-and-stdqm}\index{mathematical physics}\index{mathematicians and QM}

The previous two sections argued that the Hilbert-space framework attributed to standard quantum mechanics is doubly misnamed --- it does not distinguish QM from classical PDE continuum mechanics (\S\ref{sec:hilbert-distinction}), and the man whose name is attached to it had moved on from it within three years and built the computer architecture that his QM formalism cannot run on (\S\ref{sec:von-neumann-split}). The pattern generalises. The engagement of the mathematical community with the foundations of standard quantum mechanics has been, when one looks closely, much thinner than the formal apparatus suggests.

\paragraph{The ``one major thing then moved on'' pattern.}
The three first-rank mathematicians most often associated with the founding of the QM formalism each made one major contribution and then turned their attention elsewhere. Hilbert sponsored the Göttingen environment, gave a lecture course (winter 1926/27), and co-authored the 1928 axiomatic paper with Nordheim and von Neumann; his last major intellectual commitment was the finitist meta-mathematical programme that Gödel's 1931 incompleteness theorem then defeated. Von Neumann published the \emph{Grundlagen} in 1932, doubted the Hilbert-space framework by 1935 (\S\ref{sec:von-neumann-split}), and moved the bulk of his subsequent work into game theory, computing, and applied mathematics. Hermann Weyl published \emph{Gruppentheorie und Quantenmechanik}~\cite{Weyl1928} in 1928 --- a comprehensive group-theoretic treatment of the new mechanics --- and turned his major later work toward differential geometry, gauge theory, and the philosophy of mathematics. In each case a single substantial QM contribution stands alongside a much larger body of work in other areas, with the QM contribution itself made early in a long career and not extended.

\paragraph{The sustained mathematical-physics tradition --- real but specialised.}
A sustained tradition of mathematical work on QM does exist, mostly in the subfield that has come to be called ``mathematical physics.'' A non-exhaustive list of the principal figures:

\begin{itemize}
\item \textbf{George Mackey} (1916--2006) --- \emph{The Mathematical Foundations of Quantum Mechanics}~\cite{Mackey1963}, induced representations, group-theoretic foundations.
\item \textbf{Tosio Kato} (1917--1999) --- self-adjoint extensions of Schr\"odinger operators, perturbation theory of operators.
\item \textbf{Irving Segal} (1918--1998) --- algebraic quantum field theory, $C^*$-algebraic formulation.
\item \textbf{Arthur Wightman} (1922--2013) --- Wightman axioms for relativistic QFT; influential Princeton school.
\item \textbf{Rudolf Haag} (1922--2016) --- Haag--Kastler algebraic QFT, Haag's theorem.
\item \textbf{Michael Reed} and \textbf{Barry Simon} --- \emph{Methods of Modern Mathematical Physics}~\cite{ReedSimon} (four volumes, 1972--1978), the standard comprehensive mathematical treatment.
\item \textbf{Elliott Lieb} (1932--) --- Lieb--Thirring inequality, the proof of stability of matter, mathematical foundations of many-electron quantum mechanics.
\item \textbf{Walter Thirring} (1927--2014) --- the Lieb--Thirring co-author~\cite{LiebThirring1975}; the mathematical-physics textbook series.
\item \textbf{Edward Nelson} (1932--2014) --- stochastic mechanics as an alternative formulation of QM.
\item \textbf{Israel Gelfand} (1913--2009) --- Gelfand triples / rigged Hilbert space and the supporting operator-algebraic machinery.
\end{itemize}

This is a serious body of work and a serious tradition. But three observations sharpen what it has and has not accomplished.

\paragraph{Three observations.}

\emph{First, no first-rank mathematician made StdQM the work of their life.} The figures in the list above had QM as one substantial part of their portfolio, but rarely as their defining problem. There is no Gauss-of-number-theory or Riemann-of-geometry analogue for standard quantum mechanics --- no mathematician whose name is identified with QM the way Gauss is identified with number theory. The figure with the strongest claim is Wigner, who trained as a chemical engineer, worked in Princeton's physics department, and was awarded the 1963 Nobel Prize in \emph{Physics}; he is normally classified as a physicist with mathematical bent rather than as a mathematician.

\emph{Second, the sustained mathematical work has been on the formal infrastructure, not on the physical content.} Operator algebras, spectral theory of Schr\"odinger operators, perturbation theory, axiomatic QFT (Wightman, Haag), self-adjoint extensions (Kato) --- the volume of mathematical-physics work has been substantial and mathematically deep. But it has been largely orthogonal to the question of what the theory actually says about physical matter. The probabilistic interpretation, the role of measurement, the meaning of the wave function on $\mathbb{R}^{3N}$, the relationship to ordinary spatial intuition --- these are not problems the mathematical-physics tradition has tried to solve. They have been left, by default, to philosophers of physics (Bell, Maudlin, Albert, Bricmont).

\emph{Third, the work closest to physical content is Lieb's, and RealQM engages it directly.} The stability-of-matter programme initiated by Lieb and Thirring in 1975~\cite{LiebThirring1975} is the principal place where mathematical work on standard quantum mechanics has had bearing on a physical claim --- the claim that matter does not collapse on itself. Their proof requires antisymmetry of the many-body wave function, the Lieb--Thirring kinetic-energy bound, and a substantial technical apparatus that took years of work to assemble. RealQM's response, developed in Chapter~\ref{ch:som}, is a one-paragraph derivation of the same bound from Hardy's inequality applied to non-overlapping one-body wave functions. The Lieb--Thirring programme and the RealQM Hardy proof are doing the same physical work (ruling out collapse), in different formulations, with different technical costs. This is the one direct engagement between the mathematical-physics tradition and the present book.

\paragraph{The foundations programme: started by mathematicians, abandoned mid-stream, never picked up again.}
The serious axiomatic-mathematical foundations of standard quantum mechanics were laid down by Hilbert, von Neumann, and Weyl in a brief and intense period between 1928 and 1936 --- the Hilbert--Nordheim--von Neumann paper of 1928, Weyl's \emph{Gruppentheorie und Quantenmechanik} of 1928, von Neumann's \emph{Grundlagen} of 1932, the Birkhoff--von Neumann quantum-logic paper of 1936. That project was never completed. Von Neumann himself doubted the framework within three years of publishing it (\S\ref{sec:von-neumann-split}); Hilbert was ill from 1925 and turned his late attention to the meta-mathematical programme that G\"odel defeated in 1931; Weyl turned to differential geometry. No first-rank mathematician of comparable calibre took up the foundations of QM as their primary research problem after them. The work the mathematical-physics community has done since --- Mackey, Kato, Reed--Simon, Wightman, Haag, Lieb--Thirring --- is on the operator-theoretic infrastructure, which is mathematically deep but orthogonal to the physical-interpretational content of the theory. The interpretational and alternative-formulation work after 1936 was done by physicists --- Bohm, Everett, Bell, Zeh, GRW, 't Hooft, Rovelli --- working in a marginal subfield that the mainstream of the profession has not engaged. The combination is the institutional shape of a theory that has not been understood in a century: an abandoned mathematical-foundations programme, an infrastructure programme in place of physical content, and a marginalised physicist-foundations subfield. RealQM is offered against this background as the formulation that lets the foundations programme be picked back up --- in real 3D space, with computable equations, with the physical content of $\psi_i^2$ as charge density restored directly.

\paragraph{The pattern.}
The combination is striking when laid out together: no top-rank mathematician made StdQM their defining work; the sustained mathematical work has been on operator infrastructure rather than physical interpretation; \emph{Planck} called his own 1900 quantum hypothesis a formal assumption introduced under pressure and spent the following decade trying to derive the black-body spectrum without it (Prologue, \S\ref{sec:prologue-protagonist}); \emph{von Neumann} himself doubted the framework within three years (\S\ref{sec:von-neumann-split}); \emph{Pauli} called his own Exclusion Principle a deficiency on the occasion of receiving the Nobel Prize for it (Prologue, \S\ref{sec:prologue-strange}); \emph{Walter Kohn} declared the central object of the formalism not a legitimate scientific concept for $N \geq 10^2$--$10^3$. The description is of a discipline that has been mathematically and physically under-engaged with its own foundations for a century, propped up by a formal apparatus whose central figures --- from Planck at the start to Kohn at the end of the century --- expressed reservation almost as soon as they set down their respective contributions.

The book takes this not as a complaint about other people's work but as an opening. If the standard formulation has been under-engaged because it is structurally resistant to being engaged --- uncomputable in its central object, interpretationally opaque, with the most active subfield by deliberate choice working on the surrounding infrastructure rather than the centre --- then a formulation that can be engaged, computed, drawn on a blackboard, and varied parametrically in a browser is what the field has been missing. RealQM is offered, in this context, not as a refutation of the mathematical-physics tradition but as the formulation that lets mathematicians, physicists, chemists, and students do with quantum mechanics what they have always done with classical continuum mechanics: solve it, visualise it, modify it, teach it.

\section{Philosophy of physics: the discipline created by the gap}\label{sec:philosophy-of-physics}\index{philosophy of physics}

A further institutional fact deserves a remark, because it follows directly from the diagnosis of the previous sections. The foundational question of standard quantum mechanics --- what the wave function on $\mathbb{R}^{3N}$ actually represents, what the measurement process actually does, what it means for $|\Psi|^2$ to be a joint probability density on an abstract configuration space --- has not been taken up by either the mathematical community (which works on the operator infrastructure) or the mainstream of the physics community (which works on applications). It has, however, been taken up by a third community: the academic discipline of \emph{philosophy of physics}.

\paragraph{A discipline housed outside physics.}
Philosophy of physics is a recognised research area, with chairs, journals, conferences, and graduate programmes. A characteristic example is \textbf{Tim Maudlin} --- Professor of Philosophy at New York University, author of \emph{Philosophy of Physics: Quantum Theory} (Princeton University Press, 2019) and a substantial body of foundational work on the measurement problem, Bohmian mechanics, and the metaphysics of the wave function. Maudlin is one of the most influential current philosophers of physics on questions of quantum foundations; he holds his appointment in NYU's Department of \emph{Philosophy}, not Department of Physics. The same pattern is repeated across the discipline: David Albert (Columbia, Philosophy), David Wallace (Pittsburgh, History and Philosophy of Science), Simon Saunders (Oxford, Philosophy), Jeremy Butterfield (Cambridge, Philosophy), John Earman (Pittsburgh, History and Philosophy of Science), Hans Halvorson (Princeton, Philosophy), Jeffrey Bub (Maryland, Philosophy), Wayne Myrvold (Western Ontario, Philosophy). Their books on quantum foundations --- on the measurement problem, on the meaning of the wave function, on the interpretive choices between Copenhagen, many-worlds, Bohmian mechanics, GRW, QBism, relational QM, and their variants --- form a substantial literature with its own internal debates and citation network, published primarily in philosophy journals (\emph{Studies in History and Philosophy of Modern Physics}, \emph{Philosophy of Science}, the \emph{British Journal for the Philosophy of Science}). The notable institutional fact is where this community sits: \textbf{philosophy of physics is maintained in departments of \emph{philosophy}, not in departments of \emph{physics}}. The community that has organised itself to make sense of standard quantum mechanics is the community that does not also have to teach it as a working calculation tool.

\paragraph{Its objective and the paradox the objective contains.}
The stated objective of the discipline is to make sense --- to extract meaning, coherent logic, an interpretable physical picture --- of standard quantum mechanics. The premise of the objective is that the physics community itself has not done so: the mainstream physicist's working stance is the Mermin formula ``shut up and calculate,'' under which questions of meaning are declared either irrelevant or someone else's problem. Philosophy of physics is the discipline that has organised itself around being that someone else. Its objective is, in the literal sense, \textbf{to provide meaning, logical coherence, and interpretive structure for a physics theory that lacks all three on the physicists' own report}. Whether this is a tractable objective is open; that the discipline exists at all is itself a diagnostic.

\paragraph{What this says about the foundational situation.}
The existence of philosophy of physics as a flourishing academic subfield maintained outside physics departments is the institutional witness to the gap diagnosed in this chapter. A physical theory whose practitioners declare its physical content not understood requires an exterior community to do the work of interpretation; that community has organised, professionalised, and built its own institutions over the past several decades. Philosophy of physics did not exist as a substantial discipline before the mid-twentieth century because there was no comparable foundational gap in earlier physical theories. Maxwell's electromagnetism, Newtonian mechanics, classical thermodynamics, Einstein's relativity all admit physical interpretations that physicists themselves give without philosophical assistance. Quantum mechanics is the singular case that has spawned a parallel discipline whose specific job is to find a coherent reading of it.

\paragraph{Bohr as the founding philosopher of physics.}
The discipline did not, however, begin in philosophy departments. Its founding figure --- the figure whose work the modern philosophy of physics largely inherits and continues --- was Niels Bohr, working out of the Institute for Theoretical Physics in Copenhagen but writing on quantum mechanics in a way that was, on closer inspection, more philosophical than physical. Bohr's complementarity principle is not a physics theorem; no equation is derived from it, no calculation requires it, no measurable quantity is predicted by it. It is a philosophical thesis about the relationship between classical concepts and quantum experiments: certain pairs of classical descriptions (wave/particle, position/momentum) cannot be applied simultaneously, and the choice of which to apply is forced by the experimental arrangement. The later Bohr essays --- collected in \emph{Atomic Physics and Human Knowledge} (1958) and \emph{Essays 1958--1962 on Atomic Physics and Human Knowledge} (1963) --- are essentially philosophical writings, addressing the epistemic status of quantum descriptions, the role of the observer, the irreducibility of classical concepts in the description of measurements, and the broader implications of complementarity for biology, psychology, and culture. The Bohr--Einstein debates of the 1930s, conducted partly through the Schilpp volume, were philosophical debates between physicists, with the formalism agreed and the meaning contested.

Bohr held a physics appointment and his institute was nominally a physics institute, but the work he did on the foundations of QM from the late 1920s onward was philosophy in form and substance. The discipline of philosophy of physics --- with its measurement problem, its interpretive options, its careful logical analyses of what the formalism can and cannot say --- is, in a real sense, the academic continuation of Bohr's late-career project, carried out by professional philosophers in philosophy departments after the physicists themselves stopped doing it. \textbf{Bohr is, on this reading, the first philosopher of physics, working from inside a physics institute on a problem that the discipline now treats as belonging to philosophy.}

\paragraph{The limit of the interpretation project.}
Interpretation, however, is not reformulation. A philosophy-of-physics interpretation of standard quantum mechanics accepts the formalism --- the wave function on $\mathbb{R}^{3N}$, the Born rule, the projection postulate, the Hilbert-space machinery --- and proposes a reading of what it means. None of the interpretations on offer (Copenhagen, many-worlds, Bohmian, GRW, QBism, relational, decoherent histories) modifies the central formalism. The result is a literature in which several mutually incompatible readings coexist, each internally consistent, each defended by serious philosophers, and none accepted as definitive. This is what happens when a formalism is interpreted rather than reformulated. RealQM, by contrast, does the opposite: it changes the formalism so that the interpretive question dissolves. The wave function is a charge density on physical 3D space, in the same direct sense Maxwell's $\rho$ is; no interpretive choice is required because no probabilistic-amplitude-on-abstract-space machinery is invoked. The philosophy-of-physics community would, on this view, be working on a problem that does not arise in a formulation that stays on $\mathbb{R}^3$.

\section{Philosophy of chemistry: the same pattern on the chemistry side}\label{sec:philosophy-of-chemistry}\index{philosophy of chemistry}\index{Scerri, E.}

The institutional pattern just diagnosed for physics --- foundational and interpretive work displaced from the working discipline into a philosophy subfield --- has an exact counterpart on the chemistry side. \textbf{Philosophy of chemistry} is a recognised academic research area, with chairs, journals, conferences, and graduate programmes, housed in departments of philosophy rather than departments of chemistry. The disciplinary parallel with philosophy of physics is close enough to be worth setting out, because it confirms that the pattern is not specific to physics but reflects something more general about how the natural sciences have organised themselves around their own foundational gaps.

\paragraph{Scerri as the central figure.}
The figure most identified with the emergence of philosophy of chemistry as a distinct discipline is \textbf{Eric Scerri} --- chemist and philosopher at the University of California, Los Angeles, and the founding editor (1999) of the journal \emph{Foundations of Chemistry}. His monograph \emph{The Periodic Table: Its Story and Its Significance} (Oxford University Press, 2007) is the standard contemporary treatment of the periodic table as a philosophical-historical object; his subsequent work has developed sustained positions on the relation between chemistry and physics, the ontological status of orbitals, the role of approximation in quantum chemistry, and the question of whether chemistry reduces to physics. Scerri occupies an unusual structural position: trained as a chemist, working extensively on the philosophical foundations of chemistry, holding appointments and reaching audiences in both communities. He is, on the chemistry side, the closest analogue to the philosophy-of-physics figures (Maudlin, Albert, Wallace, Saunders) listed in the previous section.

\paragraph{The discipline's recurring questions.}
The questions that organise the philosophy-of-chemistry literature are recognisably foundational and recognisably orphaned from the working chemistry community. A non-exhaustive list:

\begin{itemize}
\item \emph{Does chemistry reduce to physics?} The standard answer in working chemistry is ``yes in principle, no in practice,'' with the in-practice clause treated as a contingent computational difficulty. The philosophy-of-chemistry literature largely rejects this reading and argues that chemistry has its own irreducible concepts --- electronegativity, oxidation state, valence, chemical bond, aromaticity, the periodic table itself --- that are not derived from but rather constrain the application of quantum-mechanical formalism to chemical systems.
\item \emph{What is an orbital?} Standard chemistry teaches orbitals as physical objects with shapes (s, p, d, f) that explain bond geometry. Standard quantum mechanics teaches that orbitals are basis functions in a Hartree--Fock or Kohn--Sham expansion --- mathematical conveniences with no direct physical referent in the many-body wave function. The disagreement is foundational, has been contested for decades, and has not been resolved by either community; the philosophy-of-chemistry literature has carried the analysis.
\item \emph{What is the chemical bond?} Lewis structures, MO theory, valence-bond theory, AIM theory (Bader), ELF, and other accounts each give a different operational answer to what a bond \emph{is}. The chemistry community treats these as alternative working tools; the philosophy-of-chemistry community treats the question itself as open.
\item \emph{Is the periodic table approximate or exact?} Whether the standard $1s, 2s, 2p, 3s, 3p, 4s, 3d \ldots$ filling order is a strict prediction of standard quantum mechanics, an approximation that fits the observed periodic table, or a phenomenological organisation that quantum mechanics fits onto but does not derive from --- the question has been pursued at length in the philosophy-of-chemistry literature and is largely absent from working chemistry curricula.
\end{itemize}

\paragraph{Structural parallel with the physics side.}
The pattern matches the philosophy-of-physics pattern almost line for line. A working community (chemists) accepts a formal framework (standard quantum chemistry, reductionist in principle, computationally arranged in practice) as adequate for getting calculations done, and treats the foundational questions about what the framework actually says as someone else's problem. A second community (philosophers of chemistry) takes those questions up, housed institutionally outside the working discipline, with its own journals and its own internal debates. The chemistry side has its founding figure as a chemist who turned philosophical (Scerri) in the same way the physics side has Bohr as a physicist who turned philosophical (Section~\ref{sec:philosophy-of-physics}). Both communities work on \emph{interpreting} an existing formalism rather than reformulating it; neither has displaced the reductionist working stance of their respective scientific communities, even though both have produced substantial literatures arguing that the reductionist stance is not actually defensible on close inspection.

\paragraph{Connection to RealQM.}
RealQM stands in an interesting place in this debate. Scerri's general position is that chemistry has its own irreducible concepts that quantum mechanics, in the standard formulation, does not derive from first principles but rather is fitted to. RealQM, by changing the formulation, returns the chemistry-relevant content of the theory to the kind of object Scerri identifies as chemically autonomous. The Level-1 RealQM model of atoms (Chapter~\ref{ch:hierarchy}) derives the radial shell structure of the periodic table from variational minimisation under spatial non-overlap, without imposing the $s/p/d/f$ orbital labels as a postulate; the closed-shell counts 2, 8, 18, 32 emerge from packing geometry (Chapter~\ref{ch:hierarchy}, §Lewis cubic atom). The Bernoulli free-boundary condition that closes RealSE (Chapter~\ref{ch:equation}) coincides geometrically with Bader's AIM zero-flux surface that the chemistry-side ontology takes seriously as the boundary of an atom in a molecule (Chapter~\ref{ch:equation}, §AIM connection). The chemical bond in RealQM is what happens when two electron territories meet at a Bernoulli boundary between two positive kernels --- a directly physical-geometric account in real 3D space, of exactly the kind the philosophy-of-chemistry critique of orbital-based StdQM has long argued for.

The book does not claim to resolve the philosophy-of-chemistry debates. It claims something narrower: that the formulation it develops makes the chemistry-side foundational questions more tractable than the StdQM formulation does, because the chemistry-relevant objects (electron territories, charge densities, free boundaries) are the primary variables rather than emergent post-processing on a $3N$-dimensional probability amplitude. If Scerri is right that chemistry has its own concepts, RealQM is a formulation that puts those concepts at the foundational level rather than at the post-processing level. Whether this counts as reductionism, anti-reductionism, or something orthogonal to the standard taxonomy is a question the philosophy-of-chemistry community is better placed to answer than this book is.

\section{Einstein and the lifelong objection}\label{sec:einstein-objection}\index{Einstein, A.}\index{EPR paradox}\index{Bohr--Einstein debates}

A short section on Einstein deserves a place in this institutional sequence because his is the most coherent and the most persistent objection to standard quantum mechanics on the historical record, and because his framing of what was wrong with the theory has organised every serious foundational discussion since.

\paragraph{Founder and critic.}
Einstein's relation to quantum mechanics is unique: he was a founder of the field (the 1905 photoelectric paper, for which he received the 1921 Nobel Prize) and then its most persistent critic from the late 1920s until his death in 1955. Unlike Schr\"odinger, whose objection was about the 3D-space picture of the wave function, Einstein's objection was structural in a different direction: it concerned the \emph{completeness} of the theory and its relation to physical reality. Quantum mechanics, on Einstein's reading, was a useful predictive instrument but not the final theoretical account of the physical world. A theory in which the most that can be said about a system is the probability distribution of measurement outcomes was, for Einstein, a placeholder for a deeper deterministic theory that the field had not yet found.

\paragraph{The Bohr--Einstein debates.}
The 5th Solvay Conference of 1927 and the 6th of 1930 contained the central public exchanges between Bohr and Einstein on the meaning of quantum mechanics --- exchanges Bohr himself later wrote up as essentially philosophical in character. Einstein deployed a sequence of thought experiments designed to show that the uncertainty relations could be violated, hence that the formalism was incomplete; Bohr each time produced a counter-analysis showing that Einstein's own gravitational redshift (the 1930 clock-in-the-box experiment) closed the apparent loophole. Bohr won the technical exchanges. But Einstein's underlying conviction --- that the probabilistic formulation is provisional, that a complete theory of the same physical situation would be deterministic and would describe what is actually happening rather than what can be measured --- was not addressed by the technical wins. He continued to hold it for the next quarter-century.

\paragraph{EPR 1935.}
The strongest formal statement of Einstein's objection is the 1935 paper with Boris Podolsky and Nathan Rosen, \emph{Can Quantum-Mechanical Description of Physical Reality Be Considered Complete?} The paper introduced what is now called an entangled state and argued, by a sufficient condition for ``elements of reality,'' that quantum mechanics either had to be incomplete or had to permit instantaneous action at a distance --- a conclusion Einstein read as a \emph{reductio} against the standard interpretation. Bohr's response, published the same year in \emph{Physical Review} under the same title, denied the EPR sufficient condition as illegitimate within the Copenhagen reading. The two papers crystallised the disagreement that subsequent foundations work would inherit: whether quantum mechanics describes physical reality directly, or only what can be said about measurement outcomes on physical reality. EPR is the founding document of quantum foundations as a discipline.

\paragraph{The Born--Einstein letters and the autobiographical notes.}
The 1916--1955 correspondence with Max Born (published 1969) is the long private record of Einstein's objection, carried out in steady form across four decades. The most famous line --- \emph{Gott w\"urfelt nicht}, ``God does not throw dice'' --- is from this correspondence and stands for the underlying realist-determinist commitment. The 1949 Schilpp volume \emph{Albert Einstein: Philosopher-Scientist}, with Einstein's ``Autobiographical Notes'' and his ``Reply to Criticisms,'' contains his mature statement: quantum mechanics is one of the great accomplishments of the human mind, and it is not what physics will look like when physics is finished.

\paragraph{Bell, 1964: Einstein's substance vindicated, his locality refuted.}
The EPR argument turned, three decades later, into the most physically substantive foundational result of the post-founders period. John Bell's 1964 inequality~\cite{BellSpeakable} showed that any local hidden-variable theory satisfying Einstein's ``elements of reality'' criterion makes statistical predictions inconsistent with those of standard quantum mechanics. The subsequent experimental tests (Aspect 1982; Aspect--Clauser--Zeilinger Nobel Prize 2022) ruled out local hidden-variable theories empirically. The technical conclusion is that Einstein's specific commitment to locality cannot be maintained alongside the empirical predictions of quantum mechanics. The substance of his broader objection --- that quantum mechanics is not the complete description of physical reality and that a deeper account remains to be found --- was not addressed by Bell's theorem and remains live. The post-Bell foundations community has split on whether Bell defeated Einstein or vindicated the seriousness of the question Einstein raised; the second reading is closer to the position this chapter takes.

\paragraph{Einstein and RealQM.}
Where does RealQM stand in this lineage? Not in the local-hidden-variable category that Bell experiments ruled out: RealQM does not propose a hidden-variable account of standard quantum mechanics. It proposes a different formulation of the underlying physics, in which the wave function is a charge density on physical 3D space rather than a probability amplitude on configuration space, and in which the probabilistic interpretation that Einstein objected to is not invoked at all. Einstein's commitment to physics as a deterministic theory of what is actually happening on physical space, rather than a probabilistic theory of what can be measured, is the commitment RealQM shares. The framework is, in this sense, the deterministic continuum reformulation Einstein would have asked for --- in the same Body-and-Soul tradition the Foreword sets out --- though carried out in 2026 with computational tools and AI collaboration he could not have anticipated, and with the experimental evidence on entanglement that he did not live to see.

\section{QED and string theory: doubling down rather than reformulating}\label{sec:qed-doubling}\index{quantum electrodynamics}\index{renormalisation}\index{string theory}

A short note on what the physics community did when standard quantum mechanics, on the foundational side, was found not to make sense. The historically observed response was not to revisit the foundations but to add another layer of theory on top --- a layer that achieved spectacular empirical accuracy and that was, in its mathematical and interpretational structure, considerably harder than the framework underneath it.

\paragraph{The QED layer.}
By the late 1920s and through the 1930s, standard quantum mechanics had been recognised as foundationally problematic --- the Bohr--Einstein debates were active, Schr\"odinger had published the cat paradox, von Neumann had privately doubted the Hilbert-space framework, Pauli had not derived his Exclusion Principle. The empirical predictions, however, were extraordinarily successful, and the physics community continued working downstream of the formalism rather than upstream of it. The next major theoretical step was quantum electrodynamics (QED), developed in the 1940s by Feynman, Schwinger, Tomonaga, and Dyson, and recognised by the 1965 Nobel Prize. QED is in many ways the empirical apex of twentieth-century physics: predictions of the electron's anomalous magnetic moment ($g{-}2$) agree with observation to twelve decimal places, the most precise theory-experiment agreement in any branch of science.

\paragraph{What the layer required.}
The mathematical and interpretational price of this empirical success was substantial. QED is built on top of the StdQM Hilbert-space formalism and adds:
\begin{itemize}
\item \emph{Quantum fields on $\mathbb{R}^4$}, replacing the wave function on $\mathbb{R}^{3N}$ with operator-valued fields acting on a Fock-space vacuum.
\item \emph{Renormalisation}: the bare predictions of the theory contain divergent integrals (ultraviolet infinities) that must be subtracted by a systematic procedure to yield finite numbers. The procedure works empirically but introduces an interpretational layer (bare charge versus dressed charge, regularisation schemes, renormalisation group) whose physical content has been debated since the 1940s.
\item \emph{Virtual particles} as intermediate states in Feynman diagrams --- non-observable entities whose status as physical processes versus calculational devices is unresolved.
\item \emph{Gauge symmetry} as the organising principle for the photon's degrees of freedom, with a constrained Hamiltonian formalism (Gupta--Bleuler, BRST quantisation) that has no direct classical analogue.
\end{itemize}
Each of these is mathematically deep and empirically necessary. None of them addresses the foundational difficulty of the underlying StdQM. Standard quantum mechanics in the 1930s did not make physical sense; QED in the 1940s did not make it make more physical sense, only enlarged what could be empirically calculated.

\paragraph{Pauli's standard, and Dirac's refusal.}
The methodological standard Pauli articulated in his 1945 Nobel Lecture --- already quoted in Chapter~\ref{ch:equation}, \S\ref{sec:pauli-nonoverlap} --- is, on inspection, an indictment of QED almost as much as it is an indictment of his own Exclusion Principle. Pauli wrote that a correct theory should not lead to infinite zero-point energies, should not use mathematical tricks to subtract infinities, and should not invent a hypothetical world that is only a mathematical fiction before formulating the correct interpretation of the actual world. Renormalisation, virtual particles, and the field-on-$\mathbb{R}^4$ vacuum picture are precisely the manoeuvres Pauli's standard rules out. Dirac, the figure whose 1929 declaration that ``the underlying physical laws of the whole of chemistry are completely known'' opens this book's foundational arc (Prologue), spent the last decades of his life refusing to accept QED on essentially Pauli's grounds. In his 1978 lectures and elsewhere, Dirac argued that subtracting infinities to obtain finite numbers is not sensible mathematics; one disregards quantities because they are small, not because they are infinitely large and one does not want them. He died in 1984 with QED accepted as orthodoxy and his own objection unanswered.

\paragraph{When QED in turn did not make sense: string theory.}
The doubling-down pattern repeated when the QED layer itself was found to be foundationally unsatisfactory. The renormalisation procedure that made QED work for electromagnetism did not extend to gravity --- attempts to quantise general relativity by analogous methods produce non-renormalisable infinities --- and the same procedure, while empirically effective for the electroweak and strong interactions, never received a clean physical interpretation. The response of the theoretical-physics community from the 1970s onward was again not to revisit the foundations but to add a further layer of theory on top. \textbf{String theory}, proposed in the early 1970s and consolidated in two successive ``string revolutions'' (1984 Green--Schwarz anomaly cancellation; 1995 Witten's M-theory and the web of dualities), replaces point particles with one-dimensional extended objects propagating in a spacetime of ten or eleven dimensions, requires supersymmetry (each known particle paired with an unobserved superpartner), and in its modern form admits something on the order of $10^{500}$ possible vacuum states --- the so-called landscape problem. Forty years of dedicated work by a substantial fraction of the world's theoretical-physics talent has not produced a single experimental prediction that distinguishes string theory from competing frameworks. The critical literature (Smolin, \emph{The Trouble with Physics}, 2006; Woit, \emph{Not Even Wrong}, 2006) has questioned whether the project as conceived can return empirical content at all. The structural pattern is the same as at the previous step: when the underlying theory was found to be foundationally unsatisfactory, the response was not to reformulate the foundations but to add a more abstract and mathematically harder layer on top.

\paragraph{The pattern continued.}
The QED and string-theory steps continue the pattern observed at every preceding stage of the chapter. A foundational difficulty in the formalism, recognised by its founders --- by Planck about the quantum hypothesis, by Pauli about Exclusion, by Schr\"odinger and Einstein about the Born interpretation, by von Neumann about Hilbert space, by Dirac about renormalisation, by Kohn about the legitimacy of $\Psi$ for $N$ large --- is left in place rather than addressed. The community builds the next layer on top, achieves new empirical successes within the layered structure (or, in the string-theory case, achieves mathematical depth without empirical content), and treats the founder's objection as a historical detail rather than as a guide to revision. QED is the most empirically successful layer; it is also the layer that most clearly violates the methodological standard the immediately previous founder (Pauli) had set out. String theory is the deepest layer yet built on top of the resulting structure; it is also the layer most distant from the empirical record. The framework developed in this book takes the opposite direction: it does not add a layer; it changes the foundation. The wave function returns to physical $\mathbb{R}^3$; the probability interpretation is removed; the antisymmetry postulate is replaced by spatial non-overlap; the configuration-space step is not taken at all. Whether the resulting framework will eventually reach the regime QED reaches (electromagnetic phenomena to twelve decimal places) is an open question; the present manuscript does not claim it. What it does claim is that the Coulomb chemistry regime --- atoms, molecules, reactions, condensed phases --- is reached without taking the steps that the founders of each successive layer regretted having taken, and without continuing the pattern of building up rather than reformulating that has organised theoretical physics for a century.

\section{The contribution restated}\label{sec:contribution-restated}

The contribution of RealQM to the chemistry-versus-physics debate is therefore concrete rather than philosophical. It is:

\begin{enumerate}
\item A specific reformulation of the $N$-electron Schr\"odinger problem (Chapter~\ref{ch:equation}) that lives on real 3D space rather than on $3N$-dimensional configuration space.
\item A demonstration (Parts~II and III) that the reformulation reproduces chemistry across atoms, bonds, hydrides, dimers, reactions, excited states, folding, and condensed phases, with accuracies of 1--10\% across the regimes tested.
\item An identification of the chemistry-specific machinery that is no longer needed: no basis sets, no exchange--correlation functionals, no fitted force fields, no transition-state machinery.
\item A working implementation that runs interactively on a laptop GPU, with the entire framework in $\sim$8000 lines of code.
\end{enumerate}

The conclusion we draw is not that the standard formulation should be abandoned but that \textbf{chemistry-as-applied-physics is achievable} when one is willing to reformulate the problem in real 3D space with an appropriate non-overlap constraint. Dirac was right; the equations were just being written in the wrong space.

% --- End of Chapter 16 ---
